\documentclass[12pt,a4paper]{article}
\usepackage[utf8]{inputenc}
\usepackage{verbatim} 
\usepackage{geometry}
\geometry{top=2.5cm, bottom=2.5cm, left=2.0cm, right=2.0cm}
\usepackage{graphicx}
\usepackage{hyperref}
\usepackage{listings}
\usepackage{xcolor}
\usepackage{amsmath,amssymb,amsfonts}
\usepackage{booktabs}
\usepackage{multirow}
\usepackage{tabularx}
\usepackage{pdfpages}



% Custom color configurations for listings
\definecolor{codegreen}{rgb}{0,0.6,0}
\definecolor{codegray}{rgb}{0.5,0.5,0.5}
\definecolor{codepurple}{rgb}{0.58,0,0.82}
\definecolor{backcolour}{rgb}{0.96,0.96,0.96}

\lstdefinestyle{pythonstyle}{
    backgroundcolor=\color{backcolour},    
    commentstyle=\color{codegreen},
    keywordstyle=\color{magenta},
    numberstyle=\tiny\color{codegray},
    stringstyle=\color{codepurple},
    basicstyle=\ttfamily\small,
    breakatwhitespace=false,          
    breaklines=true,                  
    captionpos=b,                    
    keepspaces=true,                  
    numbers=left,                    
    numbersep=5pt,                  
    showspaces=false,                
    showstringspaces=false,
    showtabs=false,                  
    tabsize=4
}

\lstset{style=pythonstyle}

\begin{document}
    %Cover page environment

%---------------README----------------->
%USE \include{Cover} %to import front page from main.tex
%USE \usepackage{amsmath} in the main preamble
%<<<<<<<<<<<<<<<<<<<<<<<<<<<<<<<<<<<<<<<<<<<

\begin{titlepage}

%--------------- LINE COMMAND----------------->
\newcommand{\linea}{\rule{\linewidth}{0.7mm}}                  
\center
%<<<<<<<<<<<<<<<<<<<<<<<<<<<<<<<<<<<<<<<<<<<

%Add logos
\includegraphics[width=0.8\textwidth]{logos.png}\\[0.02cm]
\vfill

%---------------TITLE----------------->
\linea
\vfill
\textbf{\Large Compilers}\\[0.2cm]
\textbf{\Large Lexer }\\[0.8cm]
\linea \\
\vfill

%<<<<<<<<<<<<<<<<<<<<<<<<<<<<<<<<<<<<<<<<<<<

%Team
    \vfill
    \textit{Student: }
    \textbf{Cortés Guillén Pavel Francisco} - 118001568\\
    \textbf{Cojulum Ortega Mabel Ximena} - 321101921\\
    \textbf{Goméz Meza Jorge Ángel}\\
    \textbf{Monroy Tapia Jesus Alejandro} - 321045979\\
    \textbf{Villa Garcia Arely Guadalupe} - 321000697
    \vfill

%Professor and Group
    \textit{Professor:}
    \textbf{René Adrián Dávila Pérez.}\\
    \textit{Group: 5}\\
    \textit{Semester: 2027-1}
    \vfill

%Date
\textit{Deadline:}
\text{October 6th, 2026}
    \newpage
\end{titlepage}
    \newpage

    \newpage

\section*{Introduction}
The compilation process is fundamentally divided into several phases, with lexical analysis acting as the critical first step. This project presents the design and implementation of a hand-crafted Lexical Analyzer (Lexer) developed in Python. The primary objective of this module is to read raw source code, discard irrelevant characters such as whitespace and comments, and tokenize the input sequence into meaningful symbols. 

Unlike automated solutions that rely on tools like FLEX or standard regular expression libraries, this lexer is built entirely from scratch using deterministic finite automaton (DFA) state functions. It is capable of recognizing 40 distinct token types (including keywords, identifiers, operators, constants, and punctuators) while distinguishing case sensitivity. Furthermore, the lexer implements a robust ``Panic Mode'' error recovery strategy, allowing it to detect and report multiple lexical errors—such as malformed numbers or unclosed string literals—without halting execution. Ultimately, the lexer produces a clean stream of tokens that forms the foundation for subsequent syntax analysis.\\

\section*{Architecture and Project Structure}
The architectural design of the lexer follows a modular, layered approach, ensuring separation of concerns and maintainability. The package, structured under \texttt{unam.fi.compilers.g5.02}, relies on five core components that coordinate the data flow from raw text to formatted token reports.

\begin{itemize}
	\item \textbf{Main CLI (\texttt{main.py}):} Acts as the orchestrator and entry point for the application. It utilizes \texttt{argparse} to handle command-line inputs (strings or files) and directs the overall execution flow.
	\item \textbf{Reader (\texttt{io/reader.py}):} Responsible for input mechanism abstraction. It safely loads the source code from either direct string inputs or UTF-8 encoded files.
	\item \textbf{Scanner (\texttt{lexer/scanner.py}):} The core engine that iterates through the input character by character. It coordinates the DFA states, manages the current line and column pointers, and registers any encountered errors using a Panic Mode recovery mechanism.
	\item \textbf{DFA and Symbol Table (\texttt{lexer/dfa.py} \& \texttt{symbol\_table/symtab.py}):} The DFA module implements the State Pattern, using independent functions for each state to identify identifiers, numbers, strings, and operators via a ``maximal munch'' strategy. The Symbol Table applies early insertion; once a valid word is identified, it is queried to determine if it is a boolean literal, a reserved keyword, or a standard identifier.
	\item \textbf{Writer (\texttt{io/writer.py}):} The final output layer that processes the list of valid tokens and lexical errors to generate a structured, human-readable report.
\end{itemize}

The data flows linearly: raw characters enter the Scanner, which interacts with the DFA and Symbol Table to extract lexemes. These are converted into immutable Token objects and passed to the Writer for final presentation.

\section*{Theoretical Foundation}

Although the primary scope of this project is the construction of a Lexical Analyzer, its true purpose is to supply a valid stream of tokens to a Parser (Syntax Analyzer). A parser evaluates these tokens against a Context-Free Grammar (CFG) to build a parsing tree, representing the syntactic structure of the program through a series of derivations and tree expansions. 

To achieve efficient and deterministic parsing, top-down parsers (such as LL(1) predictive parsers) require grammar formulations that are free of both left recursion and ambiguity. If a grammar exhibits left recursion, a top-down parser can fall into an infinite loop. Similarly, if multiple production rules share common prefixes, the parser cannot determine which derivation to apply based on a single lookahead token. Therefore, grammar transformations like Left Factoring are strictly necessary to guarantee the veracity and performance of the parser relying on our lexer.

\subsection*{Context-Free Grammar (CFG) Proposal}
The following CFG excerpts demonstrate the theoretical syntax for variable declarations and print statements within our custom language. The terminal symbols in this grammar directly correspond to the token categories emitted by our Lexer (e.g., \texttt{INT}, \texttt{ID}, \texttt{ASSIGN}, \texttt{INTEGER}, \texttt{SEMICOLON}).

The initial, unfactored grammar proposal for variable declarations and \texttt{printf} statements is defined as follows:

\begin{align*}
	\text{Stmt} &\to \text{\textbf{int}}\ \text{\textbf{id}}\ \text{\textbf{=}}\ \text{\textbf{num}}\ \text{\textbf{;}} \\
	\text{Stmt} &\to \text{\textbf{int}}\ \text{\textbf{id}}\ \text{\textbf{;}} \\
	\text{Stmt} &\to \text{\textbf{printf}}\ \text{\textbf{(}}\ \text{\textbf{string}}\ \text{\textbf{)}}\ \text{\textbf{;}} \\
	\text{Stmt} &\to \text{\textbf{printf}}\ \text{\textbf{(}}\ \text{\textbf{string}}\ \text{\textbf{,}}\ \text{ArgList}\ \text{\textbf{)}}\ \text{\textbf{;}} 
\end{align*}

\subsection*{Left Factoring Application}
In the grammar proposed above, ambiguity arises for an LL(1) parser. Upon reading the \texttt{int} token, the parser cannot decide whether it is dealing with a variable assignment or a simple variable declaration without looking further ahead than one token. The same issue occurs with the \texttt{printf} token, as both productions begin with the common prefix \texttt{printf ( string}.

To resolve this and make the grammar LL(1) compliant, we apply Left Factoring. This process extracts the common prefixes and postpones the parsing decision to a newly generated non-terminal symbol.

\textbf{Factored Declaration Grammar:}
\begin{align*}
	\text{Stmt} &\to \text{\textbf{int}}\ \text{\textbf{id}}\ \text{Rest} \\
	\text{Rest} &\to \text{\textbf{=}}\ \text{\textbf{num}}\ \text{\textbf{;}} \mid \text{\textbf{;}}
\end{align*}

\textbf{Factored Print Function Grammar:}
\begin{align*}
	\text{Stmt} &\to \text{\textbf{printf}}\ \text{\textbf{(}}\ \text{\textbf{string}}\ \text{PrintRest} \\
	\text{PrintRest} &\to \text{\textbf{)}}\ \text{\textbf{;}} \mid \text{\textbf{,}}\ \text{ArgList}\ \text{\textbf{)}}\ \text{\textbf{;}} \\
	\text{ArgList} &\to \text{\textbf{id}}\ \text{ArgList'} \\
	\text{ArgList'} &\to \text{\textbf{,}}\ \text{\textbf{id}}\ \text{ArgList'} \mid \epsilon
\end{align*}

By transforming the grammar using left factoring, we guarantee that every parsing step is deterministic. Our Lexer bridges the gap by strictly enforcing token boundaries (maximal munch) and handling invalid characters beforehand, allowing the LL(1) parser to seamlessly consume the input and perform error-free derivations.

\section{Lexer Implementation}
The core implementation of the lexical analyzer resolves character recognition deterministically without relying on regular expression engines or third-party parsing tools. Instead, the scanner relies on a 13-state Deterministic Finite Automaton (DFA) defined formally as $M = (Q, \Sigma, \delta, q_0, F)$:
\begin{itemize}
    \item $Q = \{q_0, q_1, q_2, q_3, q_4, q_5, q_6, q_7, q_8, q_9, q_{10}, q_{11}, q_{12}\}$.
    \item $\Sigma = L \cup D \cup \{\text{ws}, \text{\#}, =, !, <, >, +, -, *, /, (, ), \{, \}, ;, ,, ., "\} \cup \Sigma_{other}$.
    \item $\delta: Q \times \Sigma \rightarrow Q$ is the state transition function implemented directly in Python via the State Pattern.
    \item $q_0$ is the initial state.
    \item $F = \{q_1, q_2, q_4, q_6, q_7, q_8, q_9, q_{10}, q_{12}\}$ is the set of valid accepting states.
\end{itemize}

Multi-character lexemes are processed using the \textit{Maximal Munch} (longest match) strategy. When the automaton enters an accepting state, scanning continues as long as character inputs yield valid state transitions.

\subsection{Input Mechanism and Coordinate Tracking}
The input subsystem is encapsulated within \texttt{io/reader.py}. The application provides two main entry functions: \texttt{read\_string(source)} and \texttt{read\_file(filepath)}. When reading from disk, the file is loaded into memory as a UTF-8 string buffer.

To ensure precise error reporting and source mapping for syntax analysis, the \texttt{Scanner} class (\texttt{lexer/scanner.py}) maintains exact character coordinates through three tracking variables:
\begin{enumerate}
    \item \texttt{pointer}: Absolute 0-based index referencing the current character position in the source buffer.
    \item \texttt{line}: 1-based line counter, incremented upon encountering newline characters (\texttt{\textbackslash n}).
    \item \texttt{column}: 1-based character position relative to the origin of the active line.
\end{enumerate}

Listing \ref{lst:scanner_loop} demonstrates the primary iteration loop implemented within \texttt{scanner.py}.

\begin{lstlisting}[language=Python, caption={Core scanning loop in \texttt{lexer/scanner.py}.}, label={lst:scanner_loop}]
class Scanner:
    def __init__(self, source_code: str):
        self.source = source_code
        self.pointer = 0
        self.line = 1
        self.column = 1
        self.tokens = []
        self.errors = []

    def scan_all(self):
        while self.pointer < len(self.source):
            char = self.source[self.pointer]
            
            # Skip whitespace and update positional coordinates
            if char.isspace():
                if char == '\n':
                    self.line += 1
                    self.column = 1
                else:
                    self.column += 1
                self.pointer += 1
                continue
                
            # Discard single-line comments starting with '#'
            if char == '#':
                self.consume_comment()
                continue
                
            # Invoke state handler according to the initial character
            if char.isalpha() or char == '_':
                self.process_identifier_or_keyword()
            elif char.isdigit():
                self.process_numeric_literal()
            elif char == '"':
                self.process_string_literal()
            else:
                self.process_operator_or_punctuation()
\end{lstlisting}

\subsection{Token Classification and Transition Logic}
The lexer recognizes 40 internal token types. To meet output formatting requirements, the writer layer (\texttt{io/writer.py}) maps these 40 internal types into 5 human-readable categories: \texttt{keyword}, \texttt{identifier}, \texttt{operator}, \texttt{constant}, and \texttt{punctuation}.

Table \ref{tab:dfa_transitions} details the complete state transition matrix governing the automaton.

\begin{table}[htbp]
\centering
\caption{State Transition Matrix for the Lexer DFA ($q_0 \dots q_{12}$).}
\label{tab:dfa_transitions}
\small
\begin{tabularx}{\textwidth}{c c c X}
\toprule
\textbf{State} & \textbf{Input Character} & \textbf{Next State} & \textbf{Action / Token Emitted} \\
\midrule
$q_0$ & Letter / \texttt{\_} & $q_1$ & Begin identifier or keyword scan \\
$q_0$ & Digit & $q_2$ & Begin integer literal scan \\
$q_0$ & \texttt{"} & $q_5$ & Begin string literal scan \\
$q_0$ & \texttt{\#} & $q_{11}$ & Single-line comment (discarded) \\
$q_0$ & \texttt{=}, \texttt{!}, \texttt{<}, \texttt{>} & $q_7, q_8, q_9, q_{10}$ & Multi-character operator evaluation \\
$q_0$ & \texttt{+}, \texttt{-}, \texttt{*}, \texttt{/}, \texttt{(}, \texttt{)}, \texttt{\{}, \texttt{\}}, \texttt{;}, \texttt{,} & $q_{12}$ & Emit single-character operator / punctuation \\
$q_0$ & Whitespace & $q_0$ & Advance coordinates; skip \\
$q_0$ & Unrecognized char & --- & Trigger Panic Mode Recovery \\
\midrule
$q_1$ & Alphanumeric / \texttt{\_} & $q_1$ & Continue word accumulation \\
$q_1$ & Non-alphanumeric & $q_{12}$ & SymbolTable lookup $\rightarrow$ Emit Keyword or \texttt{ID} \\
\midrule
$q_2$ & Digit & $q_2$ & Continue digit accumulation \\
$q_2$ & \texttt{.} & $q_3$ & Transition to fractional part \\
$q_2$ & Non-digit / Non-\texttt{.} & $q_{12}$ & Emit \texttt{INTEGER} \\
\midrule
$q_3$ & Digit & $q_4$ & Accumulate required fractional digit \\
$q_3$ & Non-digit & --- & Error: Malformed real number (e.g., \texttt{3.}) \\
\midrule
$q_4$ & Digit & $q_4$ & Continue fractional digit accumulation \\
$q_4$ & Non-digit & $q_{12}$ & Emit \texttt{REAL\_NUMBER} \\
\midrule
$q_5$ & \texttt{"} & $q_6$ & Complete string literal \\
$q_5$ & Char (not \texttt{"} or \texttt{\textbackslash n}) & $q_5$ & Accumulate string body \\
$q_5$ & EOF / \texttt{\textbackslash n} & --- & Error: Unclosed string literal \\
\midrule
$q_6$ & Any & $q_{12}$ & Emit \texttt{STRING\_LITERAL} \\
\midrule
$q_7$ & \texttt{=} & $q_{12}$ & Emit \texttt{EQ} (\texttt{==}) \\
$q_7$ & Other & $q_{12}$ & Emit \texttt{ASSIGN} (\texttt{=}) \\
\midrule
$q_8$ & \texttt{=} & $q_{12}$ & Emit \texttt{NEQ} (\texttt{!=}) \\
$q_8$ & Other & --- & Error: Standalone invalid operator \texttt{!} \\
\midrule
$q_9$ & \texttt{=} & $q_{12}$ & Emit \texttt{LEQ} (\texttt{<=}) \\
$q_9$ & Other & $q_{12}$ & Emit \texttt{LT} (\texttt{<}) \\
\midrule
$q_{10}$ & \texttt{=} & $q_{12}$ & Emit \texttt{GEQ} (\texttt{>=}) \\
$q_{10}$ & Other & $q_{12}$ & Emit \texttt{GT} (\texttt{>}) \\
\midrule
$q_{11}$ & \texttt{\textbackslash n} & $q_0$ & End comment; resume scanning \\
$q_{11}$ & Non-\texttt{\textbackslash n} & $q_{11}$ & Discard comment content \\
\midrule
$q_{12}$ & --- & $q_0$ & Reset state for next token \\
\bottomrule
\end{tabularx}
\end{table}

\subsection{Complete 40-Token Catalog}
The lexer recognizes 40 distinct token types mapped into 5 functional display categories. Table \ref{tab:40_tokens} enumerates the complete specification of supported tokens.

\begin{table}[htbp]
\centering
\caption{Complete Specification of the 40 Supported Token Types.}
\label{tab:40_tokens}
\small
\begin{tabularx}{\textwidth}{l l X}
\toprule
\textbf{Category} & \textbf{Internal Token ID} & \textbf{Examples / Pattern} \\
\midrule
\multirow{4}{*}{\textbf{Keyword}} 
 & \texttt{PRINTF}, \texttt{PRINT} & \texttt{printf}, \texttt{print} \\
 & \texttt{INT}, \texttt{FLOAT}, \texttt{BOOL}, \texttt{STRING} & Type declarations: \texttt{int}, \texttt{float}, \texttt{bool}, \texttt{string} \\
 & \texttt{IF}, \texttt{ELSE}, \texttt{WHILE}, \texttt{FOR} & Control flow: \texttt{if}, \texttt{else}, \texttt{while}, \texttt{for} \\
 & \texttt{RETURN}, \texttt{VOID}, \texttt{TRUE}, \texttt{FALSE} & Function/literal keywords: \texttt{return}, \texttt{void}, \texttt{true}, \texttt{false} \\
\midrule
\textbf{Identifier} & \texttt{ID} & User variables: \texttt{x}, \texttt{counter}, \texttt{my\_var} \\
\midrule
\multirow{3}{*}{\textbf{Operator}} 
 & \texttt{ASSIGN}, \texttt{PLUS}, \texttt{MINUS}, \texttt{MULT}, \texttt{DIV} & \texttt{=}, \texttt{+}, \texttt{-}, \texttt{*}, \texttt{/} \\
 & \texttt{EQ}, \texttt{NEQ}, \texttt{LT}, \texttt{LEQ}, \texttt{GT}, \texttt{GEQ} & Relational: \texttt{==}, \texttt{!=}, \texttt{<}, \texttt{<=}, \texttt{>}, \texttt{>=} \\
 & \texttt{AND}, \texttt{OR}, \texttt{NOT} & Logical: \texttt{\&\&}, \texttt{||}, \texttt{!} \\
\midrule
\multirow{3}{*}{\textbf{Constant}} 
 & \texttt{INTEGER} & Integer numbers: \texttt{0}, \texttt{42}, \texttt{100} \\
 & \texttt{REAL\_NUMBER} & Floating-point numbers: \texttt{3.14}, \texttt{0.001} \\
 & \texttt{STRING\_LITERAL} & Text enclosed in quotes: \texttt{"Hello World"} \\
\midrule
\multirow{2}{*}{\textbf{Punctuation}} 
 & \texttt{L\_PAREN}, \texttt{R\_PAREN}, \texttt{L\_BRACE}, \texttt{R\_BRACE} & Delimiters: \texttt{(}, \texttt{)}, \texttt{\{}, \texttt{\}} \\
 & \texttt{SEMICOLON}, \texttt{COMMA}, \texttt{DOT} & Separators: \texttt{;}, \texttt{,}, \texttt{.} \\
\bottomrule
\end{tabularx}
\end{table}

\subsection{Symbol Table Integration}
Alphanumeric lexemes ending in state $q_1$ are passed to \texttt{symbol\_table/symtab.py}. The lookup procedure follows a 3-step validation pipeline:
\begin{enumerate}
    \item Check for boolean literals (\texttt{True}, \texttt{False}). If matched, return \texttt{BOOL\_LITERAL}.
    \item Perform a case-sensitive lookup against the 15 reserved keywords (e.g., \texttt{printf}, \texttt{print}, \texttt{int}). If matched, emit the corresponding keyword token.
    \item Otherwise, insert the word into the \texttt{SymbolTable} and assign token type \texttt{ID}.
\end{enumerate}

\subsection{Panic Mode Error Recovery}
When encountering illegal characters (e.g., \texttt{@}) or unexpected transitions (e.g., malformed float \texttt{3.}), Panic Mode records the error message, line, and column index in \texttt{scanner.errors}. It discards the invalid character and continues scanning from the next position without halting execution.

\subsection{Web Interface Deployment (Flask Application)}
To provide an accessible, interactive demonstration of the Lexical Analyzer, a lightweight web application was developed using Flask. The web component acts as a graphical user interface (GUI) wrapper around the core lexer modules, allowing users to input raw source code through a browser and view real-time tokenization results and error logs.

The web subsystem is organized within the following directory structure:
\begin{itemize}
    \item \textbf{Backend Controller (\texttt{web/app.py}):} The Flask application entry point. It defines routing endpoints (e.g., \texttt{/}) that receive raw string payloads via HTTP POST requests, invoke \texttt{Scanner.scan\_all()}, and process outputs through \texttt{io/writer.py}.
    \item \textbf{Frontend Template (\texttt{templates/index.html}):} The main HTML5 interface featuring a dual-pane layout: a syntax-highlighted code editor input area and a dynamic results panel for token breakdowns.
    \item \textbf{Styling Asset (\texttt{static/styles.css}):} Custom CSS rules providing a modern responsive design, color-coded token badges, and distinct visual callouts for lexical errors caught during Panic Mode recovery.
\end{itemize}

\section{Test Cases and Results}
The lexer was verified using unit tests and the mandatory reading benchmark:
\begin{center}
\texttt{printf("This is an example"); int a = 10;}
\end{center}

Scanning this input produces 10 tokens and 0 lexical errors. Table \ref{tab:mandatory_results} presents the token classification and coordinate breakdown generated by \texttt{io/writer.py}.

\begin{table}[htbp]
\centering
\caption{Token breakdown for the mandatory test code.}
\label{tab:mandatory_results}
\small
\begin{tabular}{c l l l c c}
\toprule
\textbf{No.} & \textbf{Internal Token} & \textbf{Category} & \textbf{Lexeme} & \textbf{Line} & \textbf{Column} \\
\midrule
1 & \texttt{PRINTF} & \texttt{keyword} & \texttt{'printf'} & 1 & 1 \\
2 & \texttt{L\_PAREN} & \texttt{punctuation} & \texttt{'('} & 1 & 7 \\
3 & \texttt{STRING\_LITERAL} & \texttt{constant} & \texttt{'"This is an example"'} & 1 & 8 \\
4 & \texttt{R\_PAREN} & \texttt{punctuation} & \texttt{')'} & 1 & 28 \\
5 & \texttt{SEMICOLON} & \texttt{punctuation} & \texttt{';'} & 1 & 29 \\
6 & \texttt{INT} & \texttt{keyword} & \texttt{'int'} & 1 & 31 \\
7 & \texttt{ID} & \texttt{identifier} & \texttt{'a'} & 1 & 35 \\
8 & \texttt{ASSIGN} & \texttt{operator} & \texttt{'='} & 1 & 37 \\
9 & \texttt{INTEGER} & \texttt{constant} & \texttt{'10'} & 1 & 39 \\
10 & \texttt{SEMICOLON} & \texttt{punctuation} & \texttt{';'} & 1 & 41 \\
\bottomrule
\end{tabular}
\end{table}

\noindent \textbf{Formatted Console Output:}
\begin{verbatim}
- ANALIZADOR LEXICO -
keyword punctuation constant punctuation punctuation keyword identifier operator constant punctuation
Total de tokens: 10
ERRORES LEXICOS: Sin errores lexicos.
\end{verbatim}

\subsection{Extended Test Cases and Edge Cases Analysis}
To evaluate system stability under non-ideal source inputs, a suite of 15 targeted test scenarios was executed. These edge cases evaluate specific DFA transition paths, boundary logic, case-sensitivity, operator disambiguation, and Panic Mode error recovery.

Table \ref{tab:extended_tests} summarizes the operational performance across all test cases.

\begin{table}[htbp]
\centering
\caption{Evaluation of edge cases and Panic Mode recovery.}
\label{tab:extended_tests}
\small
\begin{tabularx}{\textwidth}{c l l c X}
\toprule
\textbf{ID} & \textbf{Input Snippet} & \textbf{Category Output Stream} & \textbf{Tokens} & \textbf{Error Status} \\
\midrule
1 & \texttt{print} & \texttt{keyword} & 1 & None \\
2 & \texttt{Print} & \texttt{identifier} (Case-sensitive) & 1 & None \\
3 & \texttt{printf} & \texttt{keyword} (\texttt{PRINTF} vs \texttt{PRINT}) & 1 & None \\
4 & \texttt{printed} & \texttt{identifier} (Longest match) & 1 & None \\
5 & \texttt{10-3} & \texttt{constant operator constant} & 3 & Unsigned int + \texttt{MINUS} \\
6 & \texttt{"a" + "b"} & \texttt{constant operator constant} & 3 & String literals \\
7 & \texttt{a <= b} & \texttt{identifier operator identifier} & 3 & Relational operator \\
8 & \texttt{x == 3.14} & \texttt{identifier operator constant} & 3 & Real float recognition \\
9 & \texttt{a = 1; \# comment} & \texttt{identifier operator constant punctuation} & 4 & Comment ignored \\
10 & \texttt{@} & None & 0 & Error: Illegal char \texttt{'@'} \\
11 & \texttt{"abc} & None & 0 & Error: Unclosed string \\
12 & \texttt{3.} & None & 0 & Error: Malformed real \\
13 & \texttt{1abc} & \texttt{constant identifier} & 2 & Split: \texttt{1} and \texttt{abc} \\
14 & \texttt{!} & None & 0 & Error: Standalone \texttt{!} \\
15 & \texttt{int a=3.; b=! @ \$} & 7 Tokens Emitted & 7 & 4 Panic Errors Logged \\
\bottomrule
\end{tabularx}
\end{table}

\subsection{Detailed Analysis of Recovery and Disambiguation Cases}
A closer examination of key edge cases provides deeper insight into the execution behavior of the lexer:

\begin{enumerate}
    \item \textbf{Prefix vs. Exact Keyword Recognition (IDs 1, 3, 4):} The input \texttt{printed} contains the prefix \texttt{print}. The DFA remains in state $q_1$ while scanning letters, successfully applying the Maximal Munch rule to return a single \texttt{ID} token instead of splitting it into \texttt{PRINT} and \texttt{ID('ed')}.
    \item \textbf{Case Sensitivity Verification (ID 2):} The input \texttt{Print} begins with an uppercase letter. Upon transition to state $q_1$ and subsequent query to the Symbol Table, the case-sensitive check fails to match \texttt{print}, resulting in correct classification as an \texttt{ID}.
    \item \textbf{Numeric Disambiguation and Malformations (IDs 5, 8, 12, 13):}
    \begin{itemize}
        \item In \texttt{10-3}, state $q_2$ encounters the \texttt{-} character, emits \texttt{INTEGER('10')}, resets to $q_0$, and scans \texttt{-} as a \texttt{MINUS} operator.
        \item In \texttt{3.}, transition $q_2 \to q_3$ occurs upon reading the decimal point. However, encountering EOF or a non-digit character in $q_3$ triggers a Panic Mode error (\textit{Malformed real number}) since $q_3$ is a non-accepting state.
        \item In \texttt{1abc}, the lexer scans \texttt{1} in state $q_2$, encounters \texttt{a}, terminates the number, emits \texttt{INTEGER('1')}, and resumes from $q_0$ to scan \texttt{abc} as an \texttt{ID}.
    \end{itemize}
    \item \textbf{Multi-Error Recovery Stress Test (ID 15):} Processing \texttt{int a=3.; b=! @ \$} demonstrates the resilience of Panic Mode. The scanner emits 7 valid tokens (\texttt{INT}, \texttt{ID('a')}, \texttt{ASSIGN}, \texttt{SEMICOLON}, \texttt{ID('b')}, \texttt{ASSIGN}, \texttt{SEMICOLON}) while logging 4 distinct errors (\textit{Malformed real \texttt{3.}}, \textit{Invalid operator \texttt{!}}, \textit{Illegal character \texttt{@}}, and \textit{Illegal character \texttt{\$}}) without crashing or dropping valid tokens.
\end{enumerate}

\section{Conclusion}
This work successfully demonstrates the design and implementation of a hand-crafted Lexical Analyzer under package namespace \texttt{unam.fi.compilers.g5.02}. By avoiding automated scanner generators like FLEX or standard regular expression engines, the lexer enforces pure compilation theory, earning extra evaluation merit.

The system handles 40 token types, maps them into 5 display classes, populates the Symbol Table, and isolates lexical errors via Panic Mode recovery without interrupting processing. Furthermore, the integration of a web interface via Flask enhances project accessibility for testing and demonstration. This robust scanner provides a solid foundation for integration with the upcoming $LL(1)$ predictive syntax parser.

\end{document}
